% Rúbrica: Resolución del problema. Fuentes y figuras: ver el plan del documento.
\section{Arquitectura y despliegue}\label{sec:arquitectura}
% NOTA (2026-09-30), estructura del repo (desde master 4d47c3e): todo el Python vive en backend/ (api, modelado,
% analisis, referencias, despliegue, calidad, ingesta, publicacion, operacion); frontend/ es Angular;
% notebooks/ trae 00_exploracion y 01_modelo_riesgo; documento/ es este LaTeX; docs/ guarda evidencia,
% capturas, diseño e historial. Un Makefile une todo (setup, data, train, api, web, test, notebooks, doc).
% Render construye la API con Root Directory backend y despliegue/Dockerfile; Vercel publica frontend/.
% NOTA (2026-09-30), tag codigo-v3: apunta a 9b64795 y es el código que clonan los notebooks. Las salidas
% finales de los notebooks llegaron después (d1f4a65) y no están en el tag, a propósito.

\subsection{De los datos a la interfaz}
\pendiente{redacción}

\subsection{Verificaciones en el build}
\pendiente{redacción}

\subsection{Endpoints}
\pendiente{redacción}

\subsection{Costos y límites}
\pendiente{redacción}
